% ============================================================
\section{Modelo e construção das redes}\label{sec:redes}
% ============================================================

\previsto{Cerca de 1,5 página, escrita para quem nunca viu o projeto (a parte que mais
confundiu na entrega parcial). Qwen3-4B-Base e o que é cada representação: \emph{embedding} de
entrada (\cod{lex}) e saídas dos blocos 1, 18 e 36. Categorias de token (palavra inteira,
início, continuação, pontuação, número). Construção: cosseno, $k=10$ vizinhos, rede
direcionada e simetrizações por união e mútua, empates exatos da rede lexical (ocorrências do
mesmo token têm cosseno 1) e o desempate com sementes. A rede lexical de tipos sobre o
vocabulário inteiro como referência. Um diagrama pequeno da construção e o link para a página
interativa (\url{\coleta}, aba de construção). Relatório técnico, Seções 3, 5 e 6.}
